% Protocolo / solicitud para el encargado de redes.
% Compilar: pdflatex protocolo_encargado.tex
\documentclass[11pt,a4paper]{article}

\usepackage[utf8]{inputenc}
\usepackage[T1]{fontenc}
\usepackage[spanish]{babel}
\usepackage{geometry}
\usepackage{booktabs}
\usepackage{enumitem}
\usepackage{hyperref}
\usepackage{tabularx}

\geometry{margin=2.2cm}
\hypersetup{colorlinks=true,linkcolor=black,urlcolor=blue}

\title{Solicitud y protocolo de uso de tr\'afico de estudiantes\\
\large Mezcla de capturas Wireshark con ataques DDoS solo en simulaci\'on SDN}
\author{Equipo Campus-DDoS-SDN\\
Universidad Nacional del Centro del Per\'u --- Huancayo\\
Contacto: cristhianmartinezcasas@gmail.com}
\date{2026}

\begin{document}
\maketitle

\section{Objeto}
Este documento describe el uso autorizado de capturas Wireshark de
tr\'afico \textbf{benigno} de estudiantes y su mezcla con ataques DDoS
generados \textbf{\'unicamente} dentro del laboratorio simulado
(Mininet + Ryu + Docker) del proyecto Campus-DDoS-SDN.

El art\'iculo t\'ecnico del proyecto est\'a en LaTeX:
\texttt{paper/main\_es.tex} (espa\~nol) y \texttt{paper/main.tex}
(ingl\'es, Software Impacts).

\section{Qu\'e s\'i se hace}
\begin{enumerate}[leftmargin=*]
  \item Capturar tr\'afico normal de estudiantes con Wireshark en el
        segmento y horario que indique el encargado.
  \item Anonimizar el \texttt{.pcap} (remapear IP; no usar payloads
        de aplicaci\'on en el modelo).
  \item Convertir a flujos estad\'isticos y etiquetarlos
        \texttt{BENIGN}.
  \item Mezclarlos con ataques simulados
        (\texttt{synflood}, \texttt{httpflood}, \texttt{slowloris},
        \texttt{icmpflood}) contra el portal emulado
        \texttt{10.0.0.10}.
  \item Entrenar un bosque aleatorio \textbf{sin} usar la IP como
        caracter\'istica.
\end{enumerate}

\section{Qu\'e no se hace}
\begin{enumerate}[leftmargin=*]
  \item No se lanza DDoS contra la red real ni contra servidores
        institucionales.
  \item No se reenv\'ia el ataque a Internet.
  \item No se publican capturas con datos personales identificables.
\end{enumerate}

\section{Procedimiento}
\begin{enumerate}[leftmargin=*]
  \item El encargado define segmento, fecha y duraci\'on.
  \item Se entrega el archivo \texttt{.pcap}/\texttt{.pcapng} al equipo.
  \item El equipo lo coloca en \texttt{data/real/wireshark/}.
  \item Se ejecuta \texttt{bash scripts/docker/preparar\_hibrido.sh}.
  \item Se entrega el informe
        \texttt{notebooks/resultados/informe\_hibrido.txt}.
\end{enumerate}

\section{Resultados ya obtenidos (solo simulaci\'on)}
Bosque aleatorio, prueba en Ciencias Agrarias (laboratorio no visto):
F1 macro \textbf{0,979}. Se distinguen \texttt{BENIGN},
\texttt{httpflood}, \texttt{slowloris} y \texttt{synflood}.
\texttt{icmpflood} no se afirma (pocos flujos). Slowloris deja la
p\'agina sin abrir; HTTP flood la ralentiza; SYN/ICMP suben paquetes.

\section{Responsables}
\begin{tabularx}{\textwidth}{@{}lX@{}}
\toprule
Equipo & Calder\'on Leiva; Llanos Lozano; Mart\'inez Casas; Poves Mart\'inez \\
Contacto & cristhianmartinezcasas@gmail.com \\
Encargado & \textit{(completar nombre y unidad)} \\
\bottomrule
\end{tabularx}

\vspace{1.5em}
\begin{tabularx}{\textwidth}{@{}X X@{}}
\toprule
Firma del encargado / fecha & Firma del responsable del equipo / fecha \\
\vspace{2.2cm} & \\
\bottomrule
\end{tabularx}

\end{document}
